\section{Related Work}
\label{sec:related_work}

\para{Logit lens and residual stream probing.}
The logit lens~\citep{nostalgebraist2020logitlens} projects intermediate hidden states through the unembedding matrix to decode what the model ``believes'' at each layer. \citet{belrose2023tuned} address calibration failures of the raw logit lens by training per-layer affine translators (the tuned lens), achieving Spearman $\rho$ up to 0.98 with final-layer predictions while identifying ``rogue dimensions'' that dominate covariance at certain layers. \citet{yu2023neuron} decompose the \residstream into neuron-level contributions, finding that ${\sim}300$ neurons (200 attention, 100 FFN) capture 97--99\% of the predictive signal for factual knowledge, with a two-stage flow from shallow FFN to deep attention to output. \citet{li2026layer} use logit lens probing to verify genuine knowledge erasure versus surface masking, identifying three layer regimes (shallow, target, deep) with distinct erasure behaviors. Our work extends this line by tracking logit lens signals \emph{during} targeted deletion and connecting the observed information trajectory to geometric changes in the \residstream.

\para{Geometry of the residual stream.}
\citet{park2023linear} formalize the linear representation hypothesis and introduce the causal inner product $\langle \bar{\gamma}, \bar{\gamma}' \rangle_C = \bar{\gamma}^\top \Cov(\gamma)^{-1} \bar{\gamma}'$ as the correct metric for measuring concept orthogonality, showing that 26 of 27 tested concepts have linear representations. \citet{shai2024belief} demonstrate that transformers trained on HMM-generated data linearly represent belief state geometry---including fractal structures---in the \residstream, with $R^2 = 0.95$ for pairwise distance correlation. \citet{janiak2024stable} identify stable regions in activation space where small perturbations do not change model output, finding that effective editing requires pushing activations across region boundaries. We complement these geometric analyses by characterizing how PCA structure, norms, and cross-layer similarity change when specific factual information is removed.

\para{Knowledge localization and deletion.}
\citet{ghosh2024mechanistic} show that targeting the Fact Lookup (\flu) mechanism in early-to-middle MLPs (layers 2--8) for editing is far more robust than targeting output-tracing components: \flu editing limits adversarial relearning to ${\sim}20\%$ versus 47--63\% for output-tracing methods. \citet{gromov2024secretly} find that \residstream norms grow as information accumulates across layers, constraining the effective magnitude of perturbations. Our directional deletion experiments directly test whether projecting out fact directions at \flu layers is sufficient for knowledge removal, and we measure the geometric consequences that prior work leaves uncharacterized.
